\documentclass[11pt]{article}
\usepackage[margin=1in]{geometry}
\usepackage{newtxtext,newtxmath}
\usepackage{booktabs}
\usepackage{graphicx}
\usepackage{amsmath}
\usepackage[colorlinks=true,linkcolor=blue,citecolor=blue,urlcolor=blue]{hyperref}
\usepackage{caption}
\captionsetup{font=small,labelfont=bf}
\setlength{\emergencystretch}{3em}
\newcommand{\RealN}{243}
\newcommand{\RealNoisyP}{16.44}
\newcommand{\RealNoisyS}{0.107}
\newcommand{\RealBMThreeP}{24.83}
\newcommand{\RealBMThreeS}{0.559}
\newcommand{\RealNlmP}{24.78}
\newcommand{\RealNlmS}{0.551}
\newcommand{\RealTvP}{24.51}
\newcommand{\RealTvS}{0.526}
\newcommand{\RealUnetP}{29.01}
\newcommand{\RealUnetS}{0.877}
\newcommand{\RealBMThreeRawP}{21.30}
\newcommand{\RealBMThreeRawS}{0.359}
\newcommand{\RealNlmRawP}{21.28}
\newcommand{\RealNlmRawS}{0.354}
\newcommand{\RealTvRawP}{20.58}
\newcommand{\RealTvRawS}{0.333}
\newcommand{\RealGain}{4.18}
\newcommand{\RealGainRaw}{7.70}
\newcommand{\RealCI}{3.94--4.42}
\newcommand{\RealGainNlm}{4.23}
\newcommand{\RealGainTv}{4.50}
\newcommand{\RealWins}{240}
\newcommand{\RealPvalue}{2.2e-41}
\newcommand{\RealBmSigma}{0.3}
\newcommand{\RealNlmH}{0.2}
\newcommand{\RealTvW}{0.24}
\newcommand{\RealDncnnP}{28.80}
\newcommand{\RealDncnnS}{0.865}
\newcommand{\RealDncnnGain}{3.97}
\newcommand{\RealUnetGainDncnn}{0.21}
\newcommand{\RealUnetDncnnCI}{0.15--0.28}
\newcommand{\RealUnetDncnnWins}{210}
\newcommand{\LargeRealN}{20}
\newcommand{\LargeRealNoisyP}{13.35}
\newcommand{\LargeRealBMThreeP}{17.05}
\newcommand{\LargeRealNlmP}{16.89}
\newcommand{\LargeRealTvP}{16.45}
\newcommand{\LargeRealUnetP}{15.83}

\newcommand{\FmdN}{150}
\newcommand{\FmdNoisyP}{25.62}
\newcommand{\FmdNoisyS}{0.584}
\newcommand{\FmdBMThreeP}{28.07}
\newcommand{\FmdBMThreeS}{0.704}
\newcommand{\FmdNlmP}{26.96}
\newcommand{\FmdNlmS}{0.697}
\newcommand{\FmdTvP}{27.01}
\newcommand{\FmdTvS}{0.707}
\newcommand{\FmdUnetP}{31.74}
\newcommand{\FmdUnetS}{0.845}
\newcommand{\FmdGain}{3.67}
\newcommand{\FmdWins}{147}
\newcommand{\FmdPvalue}{2.9e-26}
\newcommand{\FmdCI}{3.27--4.12}
\newcommand{\FmdStrataMin}{0.4}
\newcommand{\FmdStrataMax}{11.1}
\newcommand{\FmdStrataN}{36}

\title{\bfseries Supplementary Information\\
Dose-agnostic neural deconvolution for\\
low-dose fluorescence microscopy}
\author{}
\date{\today}
\begin{document}
\maketitle

\section*{S1. Data, splits, and licenses}
\begin{table}[h]
\centering
\caption{Public datasets used in this study. BBBC014--018 low-dose measurements
are synthesized from the high-SNR references with the forward model of the main
text; BBBC039 is used for the downstream segmentation check; the AI4Life pairs
are native real noisy/clean acquisitions.}
\small
\begin{tabular}{lllll}
\toprule
Dataset & Role & FOVs used & Size (pixels) & License \\
\midrule
BBBC014 channel 1 & train/val/test & 56 / 20 / 20 & $1360\times1024$ & CC BY 3.0 \\
BBBC015 & external (5 tune / 15 test) & 20 of 144 (channel 0) & $1000\times768$ & CC BY 3.0 \\
BBBC016 (GFP) & external (5 tune / 15 test) & 20 of 72 & $512\times512$ & CC BY 3.0 \\
BBBC018 (actin, green) & external (5 tune / 15 test) & 20 of 56 & $512\times512$ & CC BY-NC-SA 3.0 \\
BBBC039 (Hoechst) & segmentation check & 100 / 50 / 50 & $520\times696$ & CC0 \\
AI4Life small pairs & real test & 243 held out & $256\times256$ & CC BY 4.0 \\
AI4Life large pairs & cross-source & 20 & $1024\times1024$ & CC BY 4.0 \\
FMD fields & second real test & 150 of 301 held out & $256\times256$ & CC BY 4.0 \\
\bottomrule
\end{tabular}
\end{table}

BBBC014, BBBC015, and BBBC016 are fluorescence images from the Broad Bioimage
Benchmark Collection; BBBC015 and BBBC016 are U2OS cytoplasm-to-nucleus assays,
and BBBC018 provides a different cell type and actin stain. BBBC015 files encode
two channels; only channel 0 is used. The BBBC018 actin signal is in the green
channel, which is extracted directly rather than averaged with the blue
channel. BBBC039 is used only for the downstream nuclei-segmentation check in
S6. The main split is at
the field-of-view level; no FOV, well, or acquisition session is shared between
training, validation, and test. Test FOVs were never used for early stopping,
baseline tuning, target selection, or performance-based figure selection; the
qualitative example in Fig.~1 uses a pre-specified test FOV.

The AI4Life-MDC25-Nuclei release (Zenodo 15624741, CC~BY~4.0) is the training
set of the 2025 AI4Life Microscopy Supervised Denoising Challenge and combines
two sources: 2{,}457
$256\times256$ noisy/clean pairs from the Jacquemet noisy-nuclei dataset and 20
$1024\times1024$ pairs derived from BBBC006. The sources differ in bit depth,
saturation, and noise model: the small references reach at most 4095 counts,
whereas the large references are clipped at 65535. We split the small pairs by
image ID into 1{,}951/243/243 training/validation/test after excluding 20 IDs
that coincide numerically with the large-field IDs as a precaution; the images
are not crops of one another (a small image correlates at most 0.33 with any
tile of the corresponding large field). Over all 1{,}951 training images, the
maximum test-train thumbnail correlation is 0.55, no pair exceeds 0.8, and the
top pair has full-resolution correlation 0.06, so there are no near duplicates
(\texttt{results/real\_leakage\_full.json}). The 20 large fields are used only
for a cross-source transfer check. The challenge's hidden test set is not
public, so the comparison uses held-out pairs from the public training release,
not a leaderboard score.

The AI4Life-MDC25-FMD release (Zenodo 15761900, CC~BY~4.0) contains 63{,}669
$256\times256$ noisy/clean pairs (multiple frames per field) from three
microscopes (CF/WF/TP), three channels and four noise levels. We extract one
frame from each of 1{,}998 fields (seed 0) and split by field, stratified by
microscope/channel/level, into 1{,}396/301/301 training/validation/test pairs.
The main comparison uses a pre-specified 150-field subset: four fields per
microscope/channel/level stratum plus six fifth fields (see
\texttt{prereg/fmd\_subset150.json}).

\section*{S2. Extended classical baselines}
The main comparison uses TV deconvolution with the known PSF; after joint
tuning of weight and ADMM parameters it is the strongest classical method at
all three doses. Table~\ref{tab:extended} reports additional baselines. The
NLM \cite{buades2005}, Wiener \cite{wiener1949}, and Richardson--Lucy
\cite{richardson1972,lucy1974} values use extended parameter grids tuned on
100 validation patches at each dose. Even with the extended grids, none of these methods
reaches the network at any dose.

\begin{table}[h]
\centering
\caption{Additional classical baselines (PSNR in dB). Doses are $P=3,30,300$
full scale. NLM, Wiener, and Richardson--Lucy use extended parameter grids
tuned on 100 validation patches; their values are archived in
\texttt{results/extended\_baselines.json} from a separate run on the same test
FOVs. They
remain below TV deconvolution and the network.}
\label{tab:extended}
\begin{tabular}{lccc}
\toprule
Method & $P=3$ & $P=30$ & $P=300$ \\
\midrule
Noisy input & 14.60 & 24.10 & 32.48 \\
TV denoising & 30.31 & 34.08 & 36.26 \\
TV deconvolution & 31.34 & 35.34 & 38.63 \\
NLM (Anscombe) & 28.55 & 33.31 & 36.41 \\
Wiener & 29.15 & 33.43 & 34.79 \\
Richardson--Lucy & 29.31 & 33.62 & 34.59 \\
U-Net (mixed-dose) & 32.27 & 36.70 & 40.01 \\
\bottomrule
\end{tabular}
\end{table}

\begin{table}[t]
\centering
\caption{Validation-selected fixed doses. The dose is the smallest one reaching
90\% success on the 20 validation FOVs and is then applied unchanged to test.
A dash under test success means the validation-selected dose did not reach 90\%
on the held-out FOVs; the factor is then also shown as a dash because it is not
claimed.}
\label{tab:protocol}
\begin{tabular}{llccccc}
\toprule
Metric & Target & TV deconv dose & U-Net dose & Factor & TV test & U-Net test \\
\midrule
PSNR & 30 dB & 10 & 3 & -- & 100\% & -- \\
PSNR & 32 dB & 30 & 10 & 3.0$\times$ & 100\% & 100\% \\
PSNR & 34 dB & 100 & 30 & 3.3$\times$ & 100\% & 100\% \\
SSIM & 0.85 & 3 & 3 & 1.0$\times$ & 95\% & 100\% \\
SSIM & 0.90 & 30 & 10 & 3.0$\times$ & 100\% & 100\% \\
SSIM & 0.92 & 30 & 30 & 1.0$\times$ & 95\% & 100\% \\
Foreground PSNR & 22 dB & 3 & 3 & -- & -- & 100\% \\
Foreground PSNR & 24 dB & 10 & 10 & 1.0$\times$ & 90\% & 100\% \\
Foreground PSNR & 26 dB & 100 & 30 & 3.3$\times$ & 100\% & 100\% \\
\bottomrule
\end{tabular}

\end{table}

\begin{table}[t]
\centering
\caption{Test-oracle protocol doses (upper bound): the smallest dose reaching
90\% success on the test FOVs. This is not a deployable selection rule and is
reported only for comparison with the validation-selected table.}
\label{tab:protocoloracle}
\begin{tabular}{llccc}
\toprule
Metric & Target & TV denoise & TV deconv & U-Net \\
\midrule
PSNR & 30 dB & 10 & 10 & 10 \\
PSNR & 32 dB & 100 & 30 & 10 \\
PSNR & 34 dB & -- & 100 & 30 \\
SSIM & 0.85 & 10 & 3 & 3 \\
SSIM & 0.90 & 30 & 30 & 10 \\
SSIM & 0.92 & 100 & 30 & 30 \\
Foreground PSNR & 22 dB & 10 & 10 & 3 \\
Foreground PSNR & 24 dB & 100 & 10 & 10 \\
Foreground PSNR & 26 dB & -- & 100 & 30 \\
\bottomrule
\end{tabular}

\end{table}

\begin{table}[t]
\centering
\caption{Ablations (PSNR in dB, matched patches). Top: original base budget,
with per-dose specialists (16k patches/60 epochs), the mixed-dose model
(24k/100), the dose-conditioned variant (24k/100), the phase-scrambled control
(24k/100, conditioned architecture), and tuned TV deconvolution. Middle: the
budget-matched specialists, mixed-dose model, dose-conditioned variant, and
blind phase-scrambled control (all 64 channels, 48k patches, 150 epochs).
Bottom: PSF width at $P=30$ for the strong mixed-dose model.}
\label{tab:ablation}
\small
\begin{tabular}{lccccc}
\toprule
Dose & Specialist (under-budget) & Mixed-dose (base) & Conditioned (base) & Scrambled & TV dec \\
\midrule
$P=3$ & 32.18 & 32.05 & 32.20 & 30.06 & 31.34 \\
$P=30$ & 36.43 & 36.27 & 36.44 & 34.03 & 35.34 \\
$P=100$ & 37.95 & 38.07 & 38.29 & 34.16 & 37.16 \\
\bottomrule
\end{tabular}
\par\vspace{8pt}
\begin{tabular}{lccccc}
\toprule
 & Specialist (strong) & Mixed-dose (strong) & Conditioned (strong) & Scrambled (strong) & TV dec \\
\midrule
$P=3$ & 32.36 & 32.27 & 32.29 & 30.19 & 31.34 \\
$P=30$ & 36.72 & 36.70 & 36.71 & 34.46 & 35.34 \\
$P=100$ & 38.68 & 38.61 & 38.66 & 36.18 & 37.16 \\
\bottomrule
\end{tabular}
\par\vspace{8pt}
\begin{tabular}{lcc}
\toprule
$\sigma$ (pixels) & U-Net & TV deconv \\
\midrule
1.5 & 36.92 & 35.85 \\
2.0 & 36.70 & 35.34 \\
2.5 & 35.76 & 34.76 \\
\bottomrule
\end{tabular}

\end{table}

The budget-matched specialists (64 channels, 48k patches, 150 epochs) reach
32.36, 36.72, and 38.68\,dB at $P=3,30,100$, within 0.1\,dB of the blind model
at the same doses (\texttt{results/spec\_strong.json}). The largest difference
is at $P=3$ (+0.09\,dB for the specialist), which is the same order as the
patch-draw variation.

\section*{S2b. Protocol selection and uncertainty}
The primary protocol selects the smallest dose reaching 90\% success on the 20
validation FOVs and applies it unchanged to the held-out FOVs
(Table~\ref{tab:protocol}). Four targets give a 3--3.3$\times$ reduction in
dose $P$ (equivalently, photon budget) with both methods meeting the test
criterion (PSNR 32 and 34\,dB, SSIM 0.90,
and foreground PSNR 26\,dB); three are ties (SSIM 0.85 and 0.92, foreground
PSNR 24\,dB); and two validation selections miss the test criterion (the U-Net
dose at PSNR 30\,dB reaches 80\%, and the TV dose at foreground PSNR 22\,dB
reaches 80\%), so those factors are not claimed. Table~\ref{tab:protocoloracle}
reports the test-oracle doses as an upper bound; because the dose grid has five
points and 90\% success corresponds to 18/20 FOVs, selection noise is
substantial. A FOV bootstrap (2,000 resamples, resampling test FOVs within each
dose) gives the oracle factor distributions in Table~\ref{tab:protocolci};
these characterize the test-oracle upper bound, not a deployable protocol
(\texttt{results/protocol\_ci.json}). No $10\times$ saving remains after the TV
baseline is tuned.

\begin{table}[h]
\centering
\caption{Test-oracle protocol doses and FOV-bootstrap factor CIs (2,000
resamples); this is an upper bound, not a validation-selected protocol.}
\label{tab:protocolci}
\small
\begin{tabular}{lcccc}
\toprule
Target & U-Net dose & TV dose & Factor [95\% CI] & Bootstrap $\geq$ threshold \\
\midrule
PSNR 30\,dB & 10 & 10 & 1.0 [1, 3.3] & 20.9\% $\geq3\times$ \\
PSNR 32\,dB & 10 & 30 & 3.0 [1, 3] & 96.6\% $\geq3\times$ \\
PSNR 34\,dB & 30 & 100 & 3.3 [3.3, 3.3] & 99.9\% $\geq3\times$ \\
SSIM 0.85 & 3 & 3 & 1.0 [1, 3.3] & 7.9\% $\geq3\times$ \\
SSIM 0.90 & 10 & 30 & 3.0 [1, 3] & 81.3\% $\geq3\times$ \\
SSIM 0.92 & 30 & 30 & 1.0 [1, 3.3] & 48.4\% $\geq3\times$ \\
FG PSNR 22\,dB & 3 & 10 & 3.3 [1, 3.3] & 79.5\% $\geq3\times$ \\
FG PSNR 24\,dB & 10 & 10 & 1.0 [1, 3] & 32.8\% $\geq3\times$ \\
FG PSNR 26\,dB & 30 & 100 & 3.3 [1, 3.3] & 97.0\% $\geq3\times$ \\
\bottomrule
\end{tabular}
\end{table}

\section*{S3. Per-FOV adaptive allocation}
An oracle with access to each FOV's full dose-response would save photons
(Table~\ref{tab:adaptive}). A ground-truth-free policy used a cheap preview,
ensemble uncertainty, and ridge or k-nearest-neighbour calibration on the 20
validation FOVs. At a 34\,dB target, the kNN policy used dose $P=24.5$ at 100\%
success versus $P=30$ for the best fixed dose at 95--96\% success, a
modest gain. At 35\,dB, the best fixed dose at 95\% test success was $P=50$,
while the six policies used $P=52$--77; no policy improved on the fixed
protocol. The preview dose is part of the final dose (the final image is
acquired by escalating from the preview), so the cost is the final dose, not
preview plus final. The gain is therefore target-dependent and not robust
enough to replace a fixed dose, which is why the main protocol uses a fixed
dose.

\begin{table}[h]
\centering
\caption{Adaptive allocation at two targets. ``Fixed'' is the dose selected on
validation for 90\% success; ``best fixed at test'' is the smallest fixed dose
that matches the reported test success; ``oracle'' uses each FOV's own curve.
Doses are $P$ in full-scale units.}
\label{tab:adaptive}
\begin{tabular}{lcccc}
\toprule
Target & Fixed & Best fixed at test & Oracle & kNN (uncertainty) \\
\midrule
PSNR 34\,dB & 50 & 30 (95--96\%) & 17 & 24.5 (100\%) \\
PSNR 35\,dB & 100 & 50 (95\%) & 29 & 68 (100\%) \\
\bottomrule
\end{tabular}
\end{table}

\section*{S4. Uncertainty details}
The ensemble contains three dose-conditioned networks trained with seeds 0, 1,
and 2 with the original 32-channel/24k-patch budget. That base-budget
conditioned model is 0.07--0.32\,dB below the strong mixed-dose main model; the
ensemble is used as a proxy for uncertainty behavior, and its absolute
calibration against the stronger model is not tested. The uncertainty score is the mean per-pixel standard deviation across
the three predictions. At $P=3$ and 10, 15 and 11 of the 20 test FOVs,
respectively, fall below 34\,dB; the corresponding AUCs are 0.96 (95\% CI
0.84--1.00) and 0.96 (0.85--1.00; \texttt{results/uncertainty\_ci.json}). At
$P=30$ only one FOV falls below the target, so the AUC is not reported. The
uncertainty also correlates with reconstruction RMSE (Spearman $\rho=0.92$ at
$P=3$). The uncertainty maps are not a proof of artifact-free reconstruction;
they identify fields where the ensemble disagrees and where additional light is
most useful.

\section*{S4b. Read noise and photon units}
At $P=30$, adding Gaussian read noise with standard deviation 0.005 or 0.01
(1.3 or 2.6 8-bit counts) changed the strong mixed-dose model from 36.70 to
36.68 and 36.66\,dB, and fair TV deconvolution from 35.34 to 35.32 and
35.30\,dB (\texttt{results/read\_noise\_fair.json}). The relative ranking is
unchanged. The dose $P$ is the Poisson scale at full 8-bit
intensity; the mean detected photons per pixel on the test patches are 0.36,
3.56, and 35.6 at $P=3,30,300$ (Poisson means; \texttt{results/photon\_units.json}).
The absolute photon counts should therefore be read as full-scale units, not as
calibrated photons per emitter.

\section*{S4c. Doses not used in training}
To test whether the model memorized the five training doses, we evaluated it on
the 20 test FOVs at $P=5,50,200$ (full-scale dose), which were never used for
training.
The measured PSNR is within 0.02\,dB of log-dose interpolation between the
neighboring trained doses (Table~\ref{tab:heldout};
\texttt{results/heldout\_doses.json}).

\begin{table}[h]
\centering
\caption{Strong mixed-dose model at held-out doses (PSNR, dB). Interpolation is
linear in $\log_{10}P$ between the two neighboring training doses.}
\label{tab:heldout}
\begin{tabular}{lcccc}
\toprule
$P$ & Neighbors & Interpolated & Measured & Residual \\
\midrule
5   & 3, 10    & 33.32 & 33.32 & $0.00$ \\
50  & 30, 100  & 37.51 & 37.50 & $-0.01$ \\
200 & 100, 300 & 39.49 & 39.47 & $-0.02$ \\
\bottomrule
\end{tabular}
\end{table}

\section*{S4d. Physical widefield PSF}
The main simulation uses a Gaussian PSF. To test a physically motivated PSF,
we repeated the $P=3/30/300$ benchmark with an in-focus scalar Airy PSF for a
circular aperture (NA 1.4, $\lambda=520$\,nm, 100\,nm pixels). TV deconvolution
was retuned on five validation FOVs with the correct Airy PSF. The
Gaussian-trained strong model loses its advantage at high dose, while a model
trained on a 50/50 Gaussian/Airy mixture (same 64-channel/48k/150 schedule)
retains positive gains at all three doses (Table~\ref{tab:airy}). On the
Gaussian benchmark the same mixed model reaches 32.26/36.60/39.83\,dB at
$P=3/30/300$, a loss of 0.01--0.18\,dB relative to the Gaussian-trained model
(\texttt{results/psf\_airy.json}, \texttt{results/psf\_airy\_mix.json},
\texttt{results/psf\_airy\_mix\_stats.json}, and
\texttt{results/blind\_psfmix\_gauss.json}).

\begin{table}[h]
\centering
\caption{Physical Airy PSF check (PSNR, dB). TV deconvolution uses the correct
Airy PSF and validation-tuned parameters. ``Gaussian-trained'' is the main
model; ``PSF-mixed'' is trained on a 50/50 Gaussian/Airy mixture.}
\label{tab:airy}
\small
\begin{tabular}{lcccc}
\toprule
$P$ & Gaussian-trained & PSF-mixed & TV (Airy) & PSF-mixed gain \\
\midrule
3   & 32.63 & 32.71 & 31.21 & +1.50 [1.35, 1.65] \\
30  & 36.64 & 37.61 & 36.23 & +1.39 [1.28, 1.48] \\
300 & 37.05 & 41.55 & 40.26 & +1.29 [1.13, 1.45] \\
\bottomrule
\end{tabular}
\end{table}

\section*{S5. Reproducibility}
\paragraph{Data.}
The four archives were downloaded from
\texttt{data.broadinstitute.org/bbbc}. SHA-256 checksums:
{\footnotesize
\begin{verbatim}
BBBC014 223c3230c164c8295f4626f82e4bd9fe394d8e9339a4f82faf697fdd9b835341
BBBC015 8c7c2abc579ff34173b49ccd06a167e1ad682f35506aa3f6e333d0336e1bad23
BBBC016 866ea3e6229a0f6fa2c9a18852b3274f7617326ecfae15d40f7f5e1123df7dd1
BBBC018 e2996b78e16de2eff74db649c1fe5335edfd82ded061e7e6afeadf27caa09df1
BBBC039 images 6f30a5d4fe38c928ded972704f085975f8dc0d65d9aa366df00e5a9d449fddd7
BBBC039 masks  f9e6043d8ca56344a4886f96a700d804d6ee982f31e2b2cd3194af2a053c2710
AI4Life Nuclei c7f2f5477a205e6020b74fded45d8f34f5383c82d261f5f401c1af58306dd7dc
AI4Life FMD    5110e958bef215e42649c2477855ea11f402842e11523ba0aebf05a634d7eb8f
\end{verbatim}
}

\paragraph{Training and evaluation.}
The main mixed-dose model was trained with
\begin{verbatim}
python src/train_unet.py --dataset bbbc014 --n-val 20 --n-test 20 \
  --ch 64 --patches 48000 --epochs 150 --val-patches 4 --mixed-doses \
  --val-all-doses --out checkpoints/bbbc014_blind_strong.pt
\end{verbatim}
The budget-matched strong conditioned variant adds \texttt{--dose-cond} and
writes \texttt{bbbc014\_dcond\_strong.pt}; budget-matched specialists use
\texttt{--photons 3}, 30, or 100 with the same 64-channel/48k/150 schedule
(\texttt{bbbc014\_spec*\_strong.pt}). The PSF-mixed variant adds
\texttt{--psf-mix} and writes \texttt{bbbc014\_blind\_psfmix.pt}. The
base-budget ablation in Table~\ref{tab:ablation} uses 24{,}000 patches and 100
epochs; its specialists use 16{,}000 patches and 60 epochs, and its
phase-scrambled control uses \texttt{--dose-cond --scramble} with the
24{,}000-patch/100-epoch schedule. The budget-matched blind control uses the
same 64-channel/48{,}000-patch/150-epoch mixed-dose schedule with
\texttt{--scramble}. Evaluation
uses \texttt{src/eval\_blind\_strong.py}, \texttt{src/eval\_dcond\_strong.py},
\texttt{src/eval\_scr\_strong.py}, \texttt{src/eval\_spec\_strong.py},
\texttt{src/eval\_matched\_ablation.py}, \texttt{src/bench\_tvdeconv.py},
\texttt{src/eval\_psf\_airy.py}, and \texttt{src/bench\_clean.py}. Figures and tables are regenerated with
\texttt{src/make\_fair\_assets.py}, \texttt{src/make\_fig1.py},
\texttt{src/make\_fig2.py}, \texttt{src/make\_fig3.py}, and
\texttt{src/make\_fig\_real.py}.

\paragraph{Hardware.}
Training and inference used an NVIDIA A100-PCIE-40GB GPU. TV deconvolution and
the CPU latency measurement used one core of the host CPU. The main split seed
is 0; the test noise seed for the matched benchmark is 2026.

\paragraph{Patch-draw sensitivity.}
The headline values use the declared ten-patch, seed-2026 benchmark. To check
the dependence on that draw, we repeated the comparison over five independent
draws of 4 or 10 patches per FOV. The paired gain remained 0.80--1.01\,dB at
$P=3$, 1.30--1.45\,dB at $P=30$, and 1.29--1.43\,dB at $P=300$; the absolute
PSNR changes with the draw, but the gain does not
(\texttt{results/patch\_draw.json}). The GPU speed factor is computed from the
rounded 1.2\,ms ($103/1.2=86$); the unrounded 1.16\,ms gives 89.

\paragraph{Code.}
The project repository contains the forward model, data loaders, splits,
baselines, training scripts, evaluation scripts, result JSON files,
checkpoints (with SHA-256 sums in \texttt{checkpoints/SHA256SUMS}), and figure
scripts; the training commands above regenerate any omitted weight. No
proprietary data are used.

\section*{S6. Downstream segmentation check}
BBBC039 provides 200 Hoechst-stained FOVs with foreground masks and an official
100/50/50 split. The restoration model was trained on 12{,}000 patches for 40
epochs with doses including $P=1$; it has no segmentation loss. TV
deconvolution and TV denoising parameters were tuned per dose on the
validation FOVs (not test), and Otsu \cite{otsu1979} was applied identically to every method.
Table~\ref{tab:seg} reports all 50 official test FOVs. With tuned TV, the
restoration model improves Dice by only about 0.02 against the best tuned TV
(0.795 versus 0.773 at $P=1$, a gain of 0.022; 0.844 versus 0.824 at $P=3$, a
gain of 0.020). This downstream check therefore supports only a small utility
gain and is not used as a main claim.

\begin{table}[h]
\centering
\caption{BBBC039 nuclei Dice at selected doses (full-scale $P$); 50 official
test FOVs, TV tuned per dose on validation, Otsu applied to every method.}
\label{tab:seg}
\begin{tabular}{lccccc}
\toprule
Method & $P=1$ & $P=3$ & $P=10$ & $P=30$ & $P=100$ \\
\midrule
Noisy + Otsu & 0.192 & 0.344 & 0.434 & 0.575 & 0.769 \\
TV denoise + Otsu & 0.772 & 0.824 & 0.857 & 0.873 & 0.882 \\
TV deconv + Otsu & 0.773 & 0.815 & 0.843 & 0.870 & 0.885 \\
U-Net + Otsu & 0.795 & 0.844 & 0.878 & 0.890 & 0.898 \\
\bottomrule
\end{tabular}
\end{table}


\section*{S7. Real noisy/clean pairs}
The real-data check uses the AI4Life-MDC25-Nuclei release
(Zenodo 15624741, CC BY 4.0). The network has the same architecture as the main
model with 48 channels in the first level and is trained on $128\times128$
patches from the 1{,}951 training pairs (16{,}000 patches, 80 epochs, batch
64, Adam, learning rate $10^{-3}$, cosine schedule, flips and $90^\circ$
rotations). Each image is normalized independently to its 0.5--99.9 percentile
range. The checkpoint is selected on the 243 validation pairs (28.6\,dB). Classical parameters are tuned on five
validation pairs; Table~\ref{tab:realtune} lists the grids and selected values
(grid means archived in \texttt{results/real\_small\_tuning\_grid.json}).
Test pairs are never used for tuning or early stopping.

\begin{table}[h]
\centering
\caption{Classical parameter tuning on five real validation pairs (mean PSNR
in dB).}
\label{tab:realtune}
\small
\begin{tabular}{lll}
\toprule
Method & Grid (PSNR, dB) & Selected \\
\midrule
BM3D $\sigma$ & 0.2:22.60, 0.3:23.30, 0.4:23.30, 0.5:23.28, 0.7:23.24, 1.0:23.20 & 0.3 \\
NLM $h$ & 0.1:22.06, 0.15:23.30, 0.2:23.35, 0.25:23.23, 0.3:23.11 & 0.2 \\
TV $w$ & 0.08:21.44, 0.12:22.12, 0.16:22.14, 0.20:22.33, 0.24:22.45, 0.32:22.45, 0.4:22.38 & 0.24 \\
\bottomrule
\end{tabular}
\end{table}

Table~\ref{tab:real} reports all \RealN\ held-out test pairs. The network
reaches \RealUnetP\,dB and \RealUnetS\ SSIM, versus \RealBMThreeP\,dB and
\RealBMThreeS\ for affine-calibrated BM3D \cite{dabov2007}, \RealNlmP\,dB and \RealNlmS\ for
NLM, \RealTvP\,dB and \RealTvS\ for TV, and \RealNoisyP\,dB and \RealNoisyS\
for the noisy input. The paired gain over BM3D is \RealGain\,dB (95\% CI
\RealCI; exact Wilcoxon $p=\RealPvalue$; \RealWins/\RealN\ wins), over NLM
\RealGainNlm\,dB, and over TV \RealGainTv\,dB. A second supervised baseline, a
17-layer DnCNN trained with the same 16{,}000 patches and 80 epochs, reaches
\RealDncnnP\,dB and \RealDncnnS\ SSIM, \RealDncnnGain\,dB above affine-calibrated
BM3D and \RealUnetGainDncnn\,dB below the U-Net (95\% CI
\RealUnetDncnnCI; \RealUnetDncnnWins/243 wins). No pair is empty (minimum
ground-truth standard deviation 0.047).

\paragraph{Affine calibration.}
The supervised networks are trained to predict the clean reference, so they can
learn a global gain and offset that BM3D, NLM, and TV cannot. For the main
comparison, each classical method receives one global affine map fitted on the
five validation pairs and applied to the held-out outputs. The raw uncalibrated
values are \RealBMThreeRawP, \RealNlmRawP, and \RealTvRawP\,dB, giving a raw
gain of about 7.7\,dB; the affine-calibrated values are \RealBMThreeP,
\RealNlmP, and \RealTvP\,dB (gain \RealGain\,dB, 95\% CI \RealCI). Both raw
and calibrated values are reported; the calibrated values are used in the main
comparison. Refitting the maps on 5, 20, 50, or all 243 validation pairs
changed the BM3D gain coefficient from 1.26 to 1.33 (NLM 1.30--1.40; TV
1.30--1.39). On a fixed 20-pair test subset, the resulting gain changed from
3.74 to 3.98\,dB (\texttt{results/real\_affine\_stability.json} and
\texttt{results/real\_affine\_243cal\_20.json}), so the five-pair choice does
not drive the reported effect.

\begin{table}[h]
\centering
\caption{Real noisy/clean pairs from the AI4Life-MDC25 Jacquemet source.
\RealN\ held-out $256\times256$ pairs; classical parameters tuned on
validation pairs. PSNR in dB.}
\label{tab:real}
\begin{tabular}{lcc}
\toprule
Method & PSNR (dB) & SSIM \\
\midrule
\multicolumn{3}{l}{\emph{AI4Life nuclei, \RealN\ held-out fields}} \\
Noisy input & \RealNoisyP & \RealNoisyS \\
TV denoising & \RealTvP & \RealTvS \\
NLM & \RealNlmP & \RealNlmS \\
BM3D & \RealBMThreeP & \RealBMThreeS \\
DnCNN (supervised) & \RealDncnnP & \RealDncnnS \\
U-Net & \RealUnetP & \RealUnetS \\
\bottomrule
\end{tabular}

\end{table}

\paragraph{Cross-source check.} Applying the same weights to the 20 large
BBBC006 fields gives \LargeRealUnetP\,dB, below BM3D (\LargeRealBMThreeP\,dB),
NLM (\LargeRealNlmP\,dB), and TV (\LargeRealTvP\,dB); the noisy input is
\LargeRealNoisyP\,dB. These fields have 16-bit saturated references and a
different noise model. The failure is a domain shift, not a contradiction of
the in-source result; it sets the scope of the real-data claim to matched
acquisition.

A label-free variant trained only on the noisy images (J-invariant block mask,
no clean targets) reached 17.89\,dB and 0.196 SSIM on the same 243 pairs,
1.45\,dB above the noisy input but 3.41\,dB below tuned BM3D; a checkerboard
mask learned the identity because the real camera noise is spatially correlated
(\texttt{results/real\_selfsup\_all.json}). The paired clean references are
therefore necessary for the present gain.

{\raggedright
\paragraph{Reproduction.} The pipeline is \texttt{src/train\_real.py},
\texttt{src/tune\_real\_classical.py}, \texttt{src/eval\_real\_small\_par.py},
\texttt{src/eval\_real\_large\_par.py}, and \texttt{src/make\_real\_assets.py};
\texttt{src/make\_fig\_real.py} regenerates Fig.~4.\par}

\section*{S7b. Second real source (FMD)}
The FMD model uses the original base budget (32 channels, 8{,}000 patches,
40 epochs), not the stronger 48-channel nuclei model, and is trained on the
1{,}396 field-split FMD
training pairs (301 validation, 301 held-out; one frame per field, and no field
crosses splits; \texttt{datasets/fmd/split.json}). Classical baselines are
tuned separately for each microscope on five validation fields. The selected
BM3D sigmas are 0.12 (CF), 0.15 (TP), and 8.0 (WF); NLM uses $h=0.1$ for all
three; TV weights are 0.08 (CF) and 0.12 (TP, WF)
(\texttt{results/fmd\_tuning\_micro.json}). On a pre-specified stratified subset
of \FmdN\ of the 301 held-out fields (four per microscope/channel/level
stratum, plus six fifth fields; the 150 fields are listed in
\texttt{prereg/fmd\_subset150.json}), the U-Net reaches \FmdUnetP\,dB and
\FmdUnetS\ SSIM, versus \FmdBMThreeP\,dB and
\FmdBMThreeS\ for BM3D, \FmdNlmP\,dB and \FmdNlmS\ for NLM, \FmdTvP\,dB and
\FmdTvS\ for TV, and \FmdNoisyP\,dB for the noisy input. The paired gain over
BM3D is \FmdGain\,dB (95\% CI \FmdCI; \FmdWins/\FmdN\ fields, Wilcoxon
$p=\FmdPvalue$; \texttt{results/fmd\_test150\_micro.json} and
\texttt{results/fmd\_stats.json}). The gain is positive in all \FmdStrataN\
strata, with group means from \FmdStrataMin\ to \FmdStrataMax\,dB; per
microscope, the mean gains are +1.87 (CF), +3.41 (TP), and +5.96 (WF), and the
largest gains are in WF. Baselines were tuned per microscope on five validation
fields; BM3D grids were 0.05--0.3 for CF/TP and 3--20 for WF, and the selected
sigmas 0.12, 0.15, and 8.0 are interior to those grids. The TV grid ended at
0.12 for TP and WF, and the selected weights lie at that boundary; we did not
extend the grid, so the secondary TV comparison may slightly understate the
tuned TV baseline. The FMD headline comparison is against BM3D. The FMD classical outputs are not affine-calibrated on FMD
validation fields, so these numbers are not directly comparable to the nuclei
table and should be read as a secondary upper
bound. The remaining 151 test fields were not evaluated with BM3D because of
its cost; the 150-field subset was pre-specified before prediction, and this is
a secondary check. FMD is a separate underlying dataset in the same AI4Life
release and required
retraining, so it is not a zero-shot cross-source transfer.

\paragraph{Reproduction.} The FMD pipeline in \texttt{src/} consists of
\texttt{fmd\_prepare.py}, \texttt{fmd\_resplit.py}, \texttt{train\_fmd.py},
\texttt{make\_fmd\_subset.py}, \texttt{tune\_fmd\_micro.py},
\texttt{eval\_fmd\_par.py}, \texttt{analyze\_fmd150.py}, and
\texttt{make\_fmd\_assets.py}.

\section*{S8. Task and dose identifiability probes}
To check whether the acquisition parameters are available to the blind network,
we trained probes to predict the dose (five levels on BBBC014) and the source
task (BBBC014/015/016/018) from a single noisy $128\times128$ patch. Each probe
uses 10 training and 10 test FOVs per class, 20 patches per FOV, with FOV-level
splits. We compare a logistic regression on 14 simple features (intensity
moments, percentiles, gradient energy, Laplacian variance, and three radial FFT
bands) with a small CNN. The dose is identified with 97.9\% (linear) and 100\%
(CNN) accuracy ($n=1000/1000$; chance 20\%), and the task with 79.5\% and
75.8\% accuracy ($n=800/800$; chance 25\%). The dose result supports the
interpretation that the mixed-dose network can infer the noise level
internally; the task result is consistent with the positive but smaller
cross-dataset transfer. Code: \texttt{src/probe\_task\_dose.py}; result:
\texttt{results/probe\_task\_dose.json}.

\bibliographystyle{unsrt}
\bibliography{references}
\end{document}
